% GENERATED FILE. Edit canonical lessons, then run npm run generate:books.
% canonical-source-hash: fnv1a64:58a09d6adf803037
% canonical-lessons: ML-C168-vitiyulla, ML-C168-itunniya, ML-C168-talnna, ML-C168-palutta, ML-C168-dhairyamulla, ML-R168-first-pass-food-animals-and-the-body, ML-R168-second-pass-clothes-things-you-use-and-describing-wo

\chapter{Wide, Narrow, Low, Ripe, Brave}
\label{ch:ml-wide-narrow-low-ripe-brave}
\hlchaptermodality{168}

\begin{chapteropening}
\textbf{By the end of this chapter:} \emph{I can say wide, narrow, low, ripe and brave.}

Complete the last lesson of chapter 168: I can say wide, narrow, low, ripe and brave.
\end{chapteropening}

\section[\texorpdfstring{\ml{വീതിയുള്ള} (vītiyuḷḷa)}{വീതിയുള്ള (vītiyuḷḷa)}]{\ml{വീതിയുള്ള} (vītiyuḷḷa) — wide}

\label{lesson:ML-C168-vitiyulla}



\begin{quote}
Before the new one: say the Malayalam for clear, then the Malayalam for deep.

On a noun, what does \emph{-āl} mean? (\textquotedblleft{}By means of,\textquotedblright{} as in \emph{kattiyāl}.)
\end{quote}

\subsection*{You'll want to know: \ml{വീതിയുള്ള}}
\textbf{\ml{വീതിയുള്ള}} — \emph{vītiyuḷḷa} — \textquotedblleft{}wide\textquotedblright{}.

\begin{grammarlens}[title={What you've built}]
One more word for this chapter.
\end{grammarlens}

\subsection*{Your turn}
\begin{itemize}
\raggedright
  \item \emph{Say it:} \emph{vītiyuḷḷa}
  \item \emph{Say it:} \emph{vītiyuḷḷa}, once more
  \item \emph{Say it:} say \emph{vītiyuḷḷa}
  \item \emph{Recall:} say the Malayalam for clear, then the Malayalam for deep, then say \emph{vītiyuḷḷa} again
\end{itemize}

\begin{tcolorbox}[breakable,colback=teal!4,colframe=teal!35!black,title={Before you move on}]
What does \ml{വീതിയുള്ള} mean? (\textquotedblleft{}Wide\textquotedblright{}.) Say it once more.
\end{tcolorbox}
\section[\texorpdfstring{\ml{ഇടുങ്ങിയ} (iṭuṅṅiya)}{ഇടുങ്ങിയ (iṭuṅṅiya)}]{\ml{ഇടുങ്ങിയ} (iṭuṅṅiya) — narrow}

\label{lesson:ML-C168-itunniya}



\begin{quote}
Before the new one: say the Malayalam for deep, then the Malayalam for wide.
\end{quote}

\subsection*{You'll want to know: \ml{ഇടുങ്ങിയ}}
\textbf{\ml{ഇടുങ്ങിയ}} — \emph{iṭuṅṅiya} — \textquotedblleft{}narrow\textquotedblright{}.

\begin{grammarlens}[title={What you've built}]
One more word for this chapter.
\end{grammarlens}

\subsection*{Your turn}
\begin{itemize}
\raggedright
  \item \emph{Say it:} \emph{iṭuṅṅiya}
  \item \emph{Say it:} \emph{iṭuṅṅiya}, once more
  \item \emph{Say it:} say \emph{iṭuṅṅiya}
  \item \emph{Recall:} say the Malayalam for deep, then the Malayalam for wide, then say \emph{iṭuṅṅiya} again
\end{itemize}

\begin{tcolorbox}[breakable,colback=teal!4,colframe=teal!35!black,title={Before you move on}]
What does \ml{ഇടുങ്ങിയ} mean? (\textquotedblleft{}Narrow\textquotedblright{}.) Say it once more.
\end{tcolorbox}
\section[\texorpdfstring{\ml{താഴ്ന്ന} (tāḻnna)}{താഴ്ന്ന (tāḻnna)}]{\ml{താഴ്ന്ന} (tāḻnna) — low}

\label{lesson:ML-C168-talnna}



\begin{quote}
Before the new one: say the Malayalam for wide, then the Malayalam for narrow.

Say twenty-one. (\emph{Irupattiyonnŭ}.)
\end{quote}

\subsection*{You'll want to know: \ml{താഴ്ന്ന}}
\textbf{\ml{താഴ്ന്ന}} — \emph{tāḻnna} — \textquotedblleft{}low\textquotedblright{}.

\begin{grammarlens}[title={What you've built}]
One more word for this chapter.
\end{grammarlens}

\subsection*{Your turn}
\begin{itemize}
\raggedright
  \item \emph{Say it:} \emph{tāḻnna}
  \item \emph{Say it:} \emph{tāḻnna}, once more
  \item \emph{Say it:} say \emph{tāḻnna}
  \item \emph{Recall:} say the Malayalam for wide, then the Malayalam for narrow, then say \emph{tāḻnna} again
\end{itemize}

\begin{tcolorbox}[breakable,colback=teal!4,colframe=teal!35!black,title={Before you move on}]
What does \ml{താഴ്ന്ന} mean? (\textquotedblleft{}Low\textquotedblright{}.) Say it once more.
\end{tcolorbox}
\section[\texorpdfstring{\ml{പഴുത്ത} (paḻutta)}{പഴുത്ത (paḻutta)}]{\ml{പഴുത്ത} (paḻutta) — ripe}

\label{lesson:ML-C168-palutta}



\begin{quote}
Before the new one: say the Malayalam for narrow, then the Malayalam for low.
\end{quote}

\subsection*{You'll want to know: \ml{പഴുത്ത}}
\textbf{\ml{പഴുത്ത}} — \emph{paḻutta} — \textquotedblleft{}ripe\textquotedblright{}.

\begin{grammarlens}[title={What you've built}]
One more word for this chapter.
\end{grammarlens}

\subsection*{Your turn}
\begin{itemize}
\raggedright
  \item \emph{Say it:} \emph{paḻutta}
  \item \emph{Say it:} \emph{paḻutta}, once more
  \item \emph{Say it:} say \emph{paḻutta}
  \item \emph{Recall:} say the Malayalam for narrow, then the Malayalam for low, then say \emph{paḻutta} again
\end{itemize}

\begin{tcolorbox}[breakable,colback=teal!4,colframe=teal!35!black,title={Before you move on}]
What does \ml{പഴുത്ത} mean? (\textquotedblleft{}Ripe\textquotedblright{}.) Say it once more.
\end{tcolorbox}
\section[\texorpdfstring{\ml{ധൈര്യമുള്ള} (dhairyamuḷḷa)}{ധൈര്യമുള്ള (dhairyamuḷḷa)}]{\ml{ധൈര്യമുള്ള} (dhairyamuḷḷa) — brave}

\label{lesson:ML-C168-dhairyamulla}



\begin{quote}
Before the new one: say the Malayalam for low, then the Malayalam for ripe.
\end{quote}

\subsection*{You'll want to know: \ml{ധൈര്യമുള്ള}}
\textbf{\ml{ധൈര്യമുള്ള}} — \emph{dhairyamuḷḷa} — \textquotedblleft{}brave\textquotedblright{}.

\begin{grammarlens}[title={What you've built}]
One more word for this chapter.
\end{grammarlens}

\subsection*{Your turn}
\begin{itemize}
\raggedright
  \item \emph{Say it:} \emph{dhairyamuḷḷa}
  \item \emph{Say it:} \emph{dhairyamuḷḷa}, once more
  \item \emph{Say it:} say \emph{dhairyamuḷḷa}
  \item \emph{Recall:} say the Malayalam for low, then the Malayalam for ripe, then say \emph{dhairyamuḷḷa} again
\end{itemize}

\begin{tcolorbox}[breakable,colback=teal!4,colframe=teal!35!black,title={Before you move on}]
What does \ml{ധൈര്യമുള്ള} mean? (\textquotedblleft{}Brave\textquotedblright{}.) Say it once more.
\end{tcolorbox}
\section[(dialogue)]{Food, animals and the body — a first pass}

\label{lesson:ML-R168-first-pass-food-animals-and-the-body}



\begin{quote}
Say the words for ripe and brave.
\end{quote}

\subsection*{Your turn}
Say these aloud:

\begin{itemize}
\raggedright
  \item a ladder, a verandah, a rainbow, a waterfall, a stone
  \item soil, a coconut palm, a banana plant, jasmine, a lotus
  \item a bear, a fox, a camel, a duck, an owl
  \item an eyebrow, a palm, rice gruel, beans, dal
\end{itemize}

\begin{tcolorbox}[breakable,colback=teal!4,colframe=teal!35!black,title={Before you move on}]
A ladder? (\textbf{\ml{കോണി}}, \emph{kōṇi}.) A verandah? (\textbf{\ml{വരാന്ത}}, \emph{varānta}.) A rainbow? (\textbf{\ml{മഴവില്ല്}}, \emph{maḻavillŭ}.) A waterfall? (\textbf{\ml{വെള്ളച്ചാട്ടം}}, \emph{veḷḷaccāṭṭaṁ}.) A stone? (\textbf{\ml{കല്ല്}}, \emph{kallŭ}.) Soil? (\textbf{\ml{മണ്ണ്}}, \emph{maṇṇŭ}.) A coconut palm? (\textbf{\ml{തെങ്ങ്}}, \emph{teṅṅŭ}.) A banana plant? (\textbf{\ml{വാഴ}}, \emph{vāḻa}.) Jasmine? (\textbf{\ml{മുല്ല}}, \emph{mulla}.) A lotus? (\textbf{\ml{താമര}}, \emph{tāmara}.) A bear? (\textbf{\ml{കരടി}}, \emph{karaṭi}.) A fox? (\textbf{\ml{കുറുക്കൻ}}, \emph{kuṟukkan}.) A camel? (\textbf{\ml{ഒട്ടകം}}, \emph{oṭṭakaṁ}.) A duck? (\textbf{\ml{താറാവ്}}, \emph{tāṟāvŭ}.) An owl? (\textbf{\ml{മൂങ്ങ}}, \emph{mūṅṅa}.) An eyebrow? (\textbf{\ml{പുരികം}}, \emph{purikaṁ}.) A palm? (\textbf{\ml{കൈവെള്ള}}, \emph{kaiveḷḷa}.) Rice gruel? (\textbf{\ml{കഞ്ഞി}}, \emph{kañji}.) Beans? (\textbf{\ml{പയറ്}}, \emph{payaṟŭ}.) Dal? (\textbf{\ml{പരിപ്പ്}}, \emph{parippŭ}.)
\end{tcolorbox}
\section[(dialogue)]{Clothes, things you use and describing words — a second pass}

\label{lesson:ML-R168-second-pass-clothes-things-you-use-and-describing-wo}



\begin{quote}
Say the words for aviyal and a pappadam.
\end{quote}

\subsection*{Your turn}
Say these aloud:

\begin{itemize}
\raggedright
  \item aviyal, a pappadam, hot, cold, long
  \item clean, sweet, bitter, spicy, sour
  \item heavy, fast, thick, soft, hard
  \item wet, dry, dark, clear, deep
  \item wide, narrow, low, ripe, brave
\end{itemize}

\begin{tcolorbox}[breakable,colback=teal!4,colframe=teal!35!black,title={Before you move on}]
Aviyal? (\textbf{\ml{അവിയൽ}}, \emph{aviyal}.) A pappadam? (\textbf{\ml{പപ്പടം}}, \emph{pappaṭaṁ}.) Hot? (\textbf{\ml{ചൂടുള്ള}}, \emph{cūṭuḷḷa}.) Cold? (\textbf{\ml{തണുപ്പുള്ള}}, \emph{taṇuppuḷḷa}.) Long? (\textbf{\ml{നീളമുള്ള}}, \emph{nīḷamuḷḷa}.) Clean? (\textbf{\ml{വൃത്തിയുള്ള}}, \emph{vṛttiyuḷḷa}.) Sweet? (\textbf{\ml{മധുരമുള്ള}}, \emph{madhuramuḷḷa}.) Bitter? (\textbf{\ml{കയ്പുള്ള}}, \emph{kaypuḷḷa}.) Spicy? (\textbf{\ml{എരിവുള്ള}}, \emph{erivuḷḷa}.) Sour? (\textbf{\ml{പുളിയുള്ള}}, \emph{puḷiyuḷḷa}.) Heavy? (\textbf{\ml{ഭാരമുള്ള}}, \emph{bhāramuḷḷa}.) Fast? (\textbf{\ml{വേഗമുള്ള}}, \emph{vēgamuḷḷa}.) Thick? (\textbf{\ml{കട്ടിയുള്ള}}, \emph{kaṭṭiyuḷḷa}.) Soft? (\textbf{\ml{മൃദുവായ}}, \emph{mṛduvāya}.) Hard? (\textbf{\ml{കഠിനമായ}}, \emph{kaṭhinamāya}.) Wet? (\textbf{\ml{നനവുള്ള}}, \emph{nanavuḷḷa}.) Dry? (\textbf{\ml{ഉണങ്ങിയ}}, \emph{uṇaṅṅiya}.) Dark? (\textbf{\ml{ഇരുണ്ട}}, \emph{iruṇṭa}.) Clear? (\textbf{\ml{തെളിഞ്ഞ}}, \emph{teḷiñña}.) Deep? (\textbf{\ml{ആഴമുള്ള}}, \emph{āḻamuḷḷa}.) Wide? (\textbf{\ml{വീതിയുള്ള}}, \emph{vītiyuḷḷa}.) Narrow? (\textbf{\ml{ഇടുങ്ങിയ}}, \emph{iṭuṅṅiya}.) Low? (\textbf{\ml{താഴ്ന്ന}}, \emph{tāḻnna}.) Ripe? (\textbf{\ml{പഴുത്ത}}, \emph{paḻutta}.) Brave? (\textbf{\ml{ധൈര്യമുള്ള}}, \emph{dhairyamuḷḷa}.)
\end{tcolorbox}
